% year: תשפ"ג
% semester: ב
% moed: א
% note: ד"ר פיטר סמובול
% source: exams/מבחנים/תשפג/מועד א תשפג סמסטר ב.pdf
% number: 1
% points: 26
% categories: func-analysis, derivatives

\begin{question}
(שאלה חובה)
\begin{enumerate}
  \item (22 נק') חקרו באופן מלא את הפונקציה: $\displaystyle f(x)=\frac{e^{2x+1}}{(x+1)^2}$ (תחום הגדרה, חיתוך עם הצירים, סימני הפונקציה, אסימפטוטות, תחומי עלייה וירידה, נקודות קיצון, תחומי קמירות ונקודות פיתול) וציירו את סקיצת הגרף של הפונקציה.
  \item (4 נק') שרטט את גרף הפונקציה $f(-|x|)$ ?
\end{enumerate}
\end{question}

\begin{hint}
בגזירה הוציאו גורם משותף: $f'(x)=\frac{e^{2x+1}\left(2(x+1)-2\right)}{(x+1)^3}$.
\end{hint}

\begin{hint}
בסעיף ב' שימו לב ש-$f(-|x|)$ היא פונקציה זוגית, ועבור $x\ge0$ היא שווה ל-$f(-x)$, כלומר שיקוף של החלק של גרף $f$ שמשמאל לציר $y$.
\end{hint}

\begin{solution}
\begin{enumerate}
  \item \textbf{תחום הגדרה:} $x\neq -1$, כלומר $D=(-\infty,-1)\cup(-1,\infty)$.

  \textbf{חיתוך עם הצירים וסימן:} $e^{2x+1}>0$ ו-$(x+1)^2>0$ בתחום, ולכן $f(x)>0$ לכל $x\in D$; אין חיתוך עם ציר $x$. חיתוך עם ציר $y$: $f(0)=e$, הנקודה $(0,e)$.

  \textbf{אסימפטוטות:}
  \begin{itemize}
    \item אנכית: $\lim_{x\to-1^\pm}f(x)=\frac{e^{-1}}{0^+}=+\infty$, ולכן $x=-1$ אסימפטוטה אנכית (משני הצדדים הפונקציה שואפת ל-$+\infty$).
    \item ב-$-\infty$: $e^{2x+1}\to 0$ ו-$(x+1)^2\to\infty$, לכן $f(x)\to0$: $y=0$ אסימפטוטה אופקית משמאל.
    \item ב-$+\infty$: לפי לופיטל (פעמיים) $\lim_{x\to\infty}\frac{e^{2x+1}}{(x+1)^2}=\lim\frac{4e^{2x+1}}{2}=\infty$, ובאופן דומה $\frac{f(x)}{x}\to\infty$, לכן אין אסימפטוטה אופקית או משופעת מימין.
  \end{itemize}

  \textbf{מונוטוניות וקיצון:} לפי כלל המנה,
  \[ f'(x)=\frac{2e^{2x+1}(x+1)^2-e^{2x+1}\cdot2(x+1)}{(x+1)^4}=\frac{2e^{2x+1}\,x}{(x+1)^3}. \]
  $f'(x)=0\iff x=0$. סימן $f'$ הוא סימן $\frac{x}{(x+1)^3}$:
  \begin{itemize}
    \item $x<-1$: מונה שלילי, מכנה שלילי $\Rightarrow f'>0$, $f$ \textbf{עולה}.
    \item $-1<x<0$: מונה שלילי, מכנה חיובי $\Rightarrow f'<0$, $f$ \textbf{יורדת}.
    \item $x>0$: $f'>0$, $f$ \textbf{עולה}.
  \end{itemize}
  לכן ב-$x=0$ יש \textbf{מינימום מקומי} $(0,e)$. אין מקסימום מקומי (ב-$x=-1$ הפונקציה אינה מוגדרת).

  \textbf{קמירות:} גזירה נוספת נותנת
  \[ f''(x)=\frac{2e^{2x+1}(2x^2+1)}{(x+1)^4}. \]
  (נגזור $f'=2e^{2x+1}\cdot x(x+1)^{-3}$: $f''=2e^{2x+1}\left[2x(x+1)^{-3}+(x+1)^{-3}-3x(x+1)^{-4}\right]=\frac{2e^{2x+1}\left[(2x+1)(x+1)-3x\right]}{(x+1)^4}$ ו-$(2x+1)(x+1)-3x=2x^2+1$.)
  מכיוון ש-$2x^2+1>0$, מתקיים $f''>0$ בכל התחום: $f$ \textbf{קמורה} על $(-\infty,-1)$ ועל $(-1,\infty)$, ו\textbf{אין נקודות פיתול}.

  \textbf{סקיצה (תיאור במילים):} משמאל הגרף יוצא מעל האסימפטוטה $y=0$ (קרוב ל-$0$ ב-$-\infty$), עולה בצורה קמורה ושואף ל-$+\infty$ כאשר $x\to-1^-$. מימין לאסימפטוטה $x=-1$ הגרף יורד מ-$+\infty$ עד המינימום $(0,e)$ ואז עולה במהירות (אקספוננציאלית) ל-$+\infty$. הגרף כולו מעל ציר $x$.

  \item נסמן $g(x)=f(-|x|)$. מתקיים $g(-x)=g(x)$, כלומר $g$ \textbf{זוגית} – הגרף סימטרי ביחס לציר $y$. עבור $x\ge0$: $g(x)=f(-x)$, כלומר הגרף של $g$ עבור $x\ge0$ הוא שיקוף ביחס לציר $y$ של גרף $f$ עבור $x\le 0$, והגרף עבור $x<0$ הוא פשוט גרף $f$ עבור $x<0$.

  בפרט: $g$ מוגדרת עבור $x\neq\pm1$, $g(0)=e$, ויש לה אסימפטוטות אנכיות $x=\pm1$ (שם $g\to+\infty$) ואסימפטוטה אופקית $y=0$ ב-$\pm\infty$. על $(-1,0)$ $g=f$ יורדת, ועל $(0,1)$ $g$ עולה; ב-$x=0$ יש מינימום מקומי $(0,e)$ (ומכיוון ש-$f'(0)=0$ אין שם "שפיץ": הנגזרות החד-צדדיות של $g$ שתיהן $0$). על $(-\infty,-1)$ $g$ עולה מ-$0$ ל-$+\infty$, ועל $(1,\infty)$ היא יורדת מ-$+\infty$ ל-$0$. $g$ קמורה בכל אחד מהקטעים.
\end{enumerate}
\end{solution}
